% ============================================================
% 个人实验报告（实验三 真机部分）——桌面物体自动分类整理（RoboMaster EP 执行线）
% 编译：xelatex 两遍（中文需 XeLaTeX；依赖 ctex、xeCJK、tikz）
% ============================================================
\documentclass[12pt,a4paper]{ctexart}

\usepackage[margin=2.4cm]{geometry}
\usepackage{amsmath,amssymb}
\usepackage{xcolor}
\usepackage{booktabs}
\usepackage{array}
\usepackage{enumitem}
\usepackage{listings}
\usepackage{hyperref}
\hypersetup{colorlinks=true,linkcolor=blue,urlcolor=blue}
\usepackage{tikz}
\usetikzlibrary{positioning,arrows.meta,shapes.geometric,fit,backgrounds}

% 让 Ⅰ(罗马数字)、≈、≥、≤ 等走 CJK 字体（lmroman 无这些字形）
\xeCJKDeclareCharClass{CJK}{"2160 -> "2160, "2248 -> "2248, "2264 -> "2265}
\emergencystretch=3em
\sloppy

\lstset{
  basicstyle=\ttfamily\small,
  frame=single,
  breaklines=true,
  columns=fullflexible,
  numbers=left,
  numberstyle=\tiny\color{gray},
  aboveskip=4pt,belowskip=2pt,
}

\tikzset{
  proc/.style={draw, rounded corners=2pt, fill=blue!6, align=center, font=\small, inner sep=4pt},
  sys/.style={draw, rounded corners=3pt, fill=green!7, thick, align=center, font=\small, inner sep=5pt},
  seam/.style={draw, rounded corners=2pt, fill=orange!14, thick, dashed, align=center, font=\small, inner sep=4pt},
  hw/.style={draw, rounded corners=3pt, fill=red!7, thick, align=center, font=\small, inner sep=5pt},
  dec/.style={draw, diamond, aspect=2.2, fill=orange!10, align=center, font=\scriptsize, inner sep=2pt},
  arr/.style={-{Stealth[length=2.6mm]}, thick},
  lbl/.style={font=\scriptsize\itshape},
  grp/.style={draw, dashed, rounded corners=4pt, inner sep=6pt},
}

\begin{document}

% ---------------- 封面信息 ----------------
\begin{center}
  {\LARGE\bfseries 个人实验报告（实验三）}\\[6pt]
  {\Large 桌面物体自动分类整理 —— 真机执行（RoboMaster EP）}
\end{center}

\begin{center}
\renewcommand{\arraystretch}{1.35}
\begin{tabular}{@{}l p{3.5cm} l p{3.5cm} l p{3.5cm}@{}}
  班级 & \underline{2024届中外2班} & 姓名 & \underline{马晨曦} & 学号 & \underline{24020036047} \\
  小组角色 & \multicolumn{3}{l}{任务控制与系统集成 · \textbf{真机执行主责}} & 日期 & \underline{2026-09-24} \\
\end{tabular}
\end{center}

\section{实验的基本要求}

本实验为机器人集成小组项目Ⅰ《桌面物体自动分类整理》，要求在实验一（视觉检测模型）与实验二
（机械臂定点抓取）的基础上，做出一个\textbf{“看得见—判得准—抓得起—放得对”的全自动整理系统}：
相机看桌面，识别出网格里摆放的多类物体，机器人依次把每个物体按类别放进对应料盒，直到桌面清空。

实验要求分为\textbf{阶段一 仿真开发}与\textbf{阶段二 真机验证}两段，且
\textbf{阶段二的准入条件 = 阶段一仿真验收通过}（顺序不可颠倒）。\textbf{本报告只写阶段二真机部分}；
阶段一仿真的验收结论只在下一小节照录一份，作为准入条件的凭据，细节归属仿真线负责人。

\textbf{本版范围声明：}阶段一仿真已于 2026-09-24 验收通过（见下），本报告据此成立。
阶段二的\textbf{现场实测数据}——\texttt{config\_real/} 的坐标与位姿表、
现场运行产生的 \texttt{real\_run.json} 与 \texttt{records.jsonl}、
分类结果记录表——\textbf{须于现场作业后补录，本版不含，亦不预填任何估计值}；
未标定时真机入口拒绝启动（§3.4），因此不存在“拿编出来的坐标跑一遍”的可能。

\subsection*{阶段一仿真的验收结果（阶段二的准入条件）}

仿真线（mechArm + Gazebo，ROS\,2）的\textbf{记分轮}已于 2026-09-24 跑通，
\begin{center}
\begin{tabular}{@{}ccccc@{}}
\toprule
格号 & 类别 & 目标料盒 & 落料槽位 & 结果 \\
\midrule
c1 & cup   & bin\_cup   & 0 & 成功 \\
c2 & mouse & bin\_mouse & 0 & 成功 \\
c3 & cup   & bin\_cup   & 1 & 成功 \\
c4 & cup   & bin\_cup   & 2（第 2 层） & 成功 \\
c5 & mouse & bin\_mouse & 1（叠层） & 成功 \\
c6 & cup   & bin\_cup   & 3 & 失败（放置偏差 52/87/17\,mm） \\
\bottomrule
\end{tabular}
\end{center}
结果文件 \texttt{exp3\_logs/run\_20260924\_112831/result.json} 给出
\texttt{verdict=PASS}：场上物体 6 个，正确放入对应料盒 5 个（及格线 5），
\textbf{规划失败 0 次}。

\textbf{异常轮}（\texttt{run\_20260924\_114058}）单独跑：c3 一格故意空置，
\texttt{verdict=TRIAL}。空网格被跳过的证据不是某一行写了什么，而是
\texttt{records.jsonl} 里\textbf{没有 c3 这一行}、\texttt{result.json} 的
\texttt{reasons} 里也没有 c3 的键；其余 5 个全部放置成功。这一轮必然是 TRIAL 而非
PASS —— 判定规则只对\textbf{完整}的验收轮下结论，所以异常轮必须与记分轮分开跑。

\textbf{口径说明（必须分开陈述）。}本次仿真跑在没有 GPU 的云机容器里，Gazebo 走软件渲染，
跑到第 3$\sim$4 轮会出现\textbf{画面停滞}：帧戳照常推进而画面内容一字不变，依赖图像的
HSV 检测因此失效。为把故障隔离在渲染这一段、继续验证下游链路，记分轮改用\textbf{真值投影}
生成检测框：物块世界坐标经同一份已标定的相机参数投影为像素，再走完全相同的
“像素 $\to$ 桌面 $\to$ 格子”映射，相机标定与映射链路未作任何改动。因此这一轮
\textbf{能证明}的是相机标定、投影与网格映射正确，以及规划—取放—落料—回航这条执行链可行；
\textbf{不能}声称 HSV 检测算法在整轮任务中稳定有效 —— 渲染健康的第 1 轮 HSV 实测
6/6 正确映射，这份证据单独保留。

\textbf{遗留：}叠第 2 层时物块之间发生互穿，LCP 求解器迭代数冲到 122847 并把已放好的物块弹飞；
第 2 层落差 \texttt{dz}=0.036\,m 仅比杯高 0.035\,m 多 1\,mm，加大落差或让层间 xy 错开是下一步
要改的地方。c6 是记分轮唯一失败项。

\subsection*{真机任务（docx §五）}

\begin{itemize}[leftmargin=1.6em]
  \item 在\textbf{真实作业区}（经确认在实验室地面上进行）标出取物网格与不少于 2 个分类料盒，
        摆放 \textbf{2 类共 6 个实物}；
  \item 使用\textbf{前序实验训练的检测模型}识别物体类别与所在网格；
  \item 使用\textbf{真机抓取程序}，把不同类别放入对应料盒。
\end{itemize}

\subsection*{真机验收标准（docx §六）}

\begin{itemize}[leftmargin=1.6em]
  \item 6 个物体中\textbf{至少 5 个被正确分类整理}；
  \item \textbf{正常任务过程中不得进行人工类别判断和动作触发}；
  \item 正确处理\textbf{至少 2 类异常}（空网格、未识别物体应被跳过；抓取失败或目标不可达时
        应返回安全位置或停止）；
  \item 提供一个入口启动图像、识别、机械臂和任务控制；
  \item 全程保存\textbf{识别、抓取和任务状态日志}。
\end{itemize}

\subsection*{真机平台的确定}

docx §实施原则写的“真实相机、Jetson 和 mechArm”是 mechArm 时代的条款。本组真机平台
\textbf{沿用实验二已完成 C5 验收的 RoboMaster EP 工程形态}（教师已于实验二阶段确认换平台，
本实验沿用同一口径）：

\begin{itemize}[leftmargin=1.6em]
  \item EP 为\textbf{工程形态}：全向底盘 + 二轴机械臂 + 夹爪，由\textbf{大疆官方 Python SDK}
        直连控制，\textbf{真机侧不使用 ROS\,2}，因此也不需要 Jetson；
  \item 相机为\textbf{USB 接口摄像头}（与实验一同一款），固定于工作区上方；
  \item 仿真线仍为 mechArm + Gazebo（ROS\,2），由组员甲负责。
\end{itemize}

\textbf{注：}本人承担\textbf{任务控制（状态机）/ 接口契约 / 日志与验收口径}这条主线，
并任\textbf{真机执行主责}——EP 侧标定、执行后端（\texttt{EPPickExecutor}）、真机入口
（\texttt{run\_real.py}）与现场运行录制。视觉模型来自实验一公开仓库，相机识别与仿真场景
分别由两位组员完成。\textbf{本报告围绕真机执行这条线展开。}

\section{个人分工}

本人在真机阶段负责以下五件事：

\begin{enumerate}[leftmargin=1.8em]
  \item \textbf{EP 执行后端 \texttt{EPPickExecutor}}：实现执行契约的\textbf{六段动作}
        （下探/夹取/抬离/搬运/松放/退回），每一段都落到实验二已验证过的 EP 原语上，
        并在失败时按段号映射为动作级原因；
  \item \textbf{EP 侧标定与位姿表 \texttt{config\_real/}}：底盘里程计位姿（6 网格 + 2 料盒 +
        归位点）、臂的两档高度、夹爪档位，全部由现场实测回填；
  \item \textbf{真机入口 \texttt{run\_real.py}}：把“相机 $\to$ 检测 $\to$ 任务控制 $\to$ EP”
        串成一条命令，含启动前的\textbf{标定闸门}与异常收尾；
  \item \textbf{未标定拒跑闸门}：把“点位到底标没标过”从一句口头约定变成\textbf{可执行的
        启动前检查}，避免拿没人量过的坐标去驱动真机械臂；
  \item \textbf{真机运行与录制}：现场标定、单周期人盯跑通、全自动一整轮录制，
        以及验收日志与结果记录。
\end{enumerate}

\textbf{与本人的仿真线工作的关系：}任务控制核心（\texttt{sort\_core}）、接口契约
（\texttt{PickPlace}）、原因枚举与日志 schema 是我在阶段一建立的，两条线\textbf{共用同一份}；
真机阶段我做的是\textbf{把执行侧换成 EP}，核心一行未改（见 §3.1）。

组内分工概览（真机阶段）：\textbf{本人} = 真机执行主责（EP 标定 / 执行后端 / 入口 / 运行录制）
+ 任务控制与契约；\textbf{组员甲} = 仿真场景搭建与联调、仿真演示视频、仿真与真机的尺寸对齐；
\textbf{组员乙} = 识别链路（真模型）、真机相机标定、分类结果记录表。

\section{项目实现的原理}

\subsection{双线架构：一套任务核心，两个注入缝隙}

真机不是另写一套系统。任务控制核心 \texttt{TaskController} 只依赖两个函数缝隙
\texttt{scan()}（报一帧桌上检测）与 \texttt{pick\_place(cell\_id, cls)}（把指定格子的物体
取走放好），\textbf{自己不 import 任何 ROS 类型}。真机线要做的只是\textbf{换掉这两个缝隙的
实现}（图~\ref{fig:arch}）：

\begin{itemize}[leftmargin=1.6em]
  \item 缝隙 1：由“订阅 \texttt{Detection2DArray}”改为\textbf{直接从 USB 相机取帧并调检测模型}；
  \item 缝隙 2：由“发 \texttt{PickPlace} Action”改为\textbf{直接调用 \texttt{EPPickExecutor}}。
\end{itemize}

这条“决策与通信解耦”的设计是阶段一定的，真机阶段兑现了它的价值：\textbf{契约、原因枚举、
日志 schema 一行未改}，换平台只换了接线。

\begin{figure}[htbp]
  \centering
  \begin{tikzpicture}[node distance=3.4mm]
    \node[sys, text width=11.4cm] (cam) {\textbf{相机（USB，固定于工作区上方）}\\
      \texttt{cv2} 取帧（640$\times$480）$\rightarrow$ \textbf{检测模型}（cup / mouse）
      $\rightarrow$ 检测列表};
    \node[seam, text width=11.4cm, below=of cam] (gap1) {\textbf{缝隙 1：\texttt{scan()}}
      —— \texttt{camera.make\_scan}（框中心 $\to$ 网格，或直接取检测列表）};
    \node[sys, text width=11.4cm, below=of gap1] (ctl) {\textbf{任务控制核心 \texttt{TaskController}}
      （纯 Python，无 ROS 依赖）\\ 扫描 $\rightarrow$ 逐条判定 $\rightarrow$ 选目标 $\rightarrow$
      执行 $\rightarrow$ 循环，直到清空 / 异常退出};
    \node[seam, text width=11.4cm, below=of ctl] (gap2) {\textbf{缝隙 2：\texttt{pick\_place(cell\_id, cls)}}
      —— \texttt{EPPickPlace}（goal 只给格号与类别）};
    \node[sys, fill=blue!8, text width=11.4cm, below=of gap2] (exe) {\textbf{\texttt{EPPickExecutor}}
      ：把格号/类别解成\textbf{底盘位姿}，执行六段动作 $\rightarrow$ \textbf{EP 原语}
      （\texttt{goto / arm\_moveto / gripper\_*}）};
    \node[hw, text width=11.4cm, below=of exe] (rm) {\textbf{RoboMaster EP 工程形态}：全向底盘
      + 二轴臂 + 夹爪\quad$\vert$\quad\textbf{大疆 Python SDK}（\textbf{无 ROS\,2}）};
    \node[sys, fill=green!10, text width=11.4cm, below=of rm] (log) {\textbf{日志/验收}：
      \texttt{exp3\_logs\_real/run\_*/}\ \texttt{records.jsonl} $+$ \texttt{result.json}
      $+$ \texttt{real\_run.json}};

    \draw[arr] (cam)--(gap1); \draw[arr] (gap1)--(ctl); \draw[arr] (ctl)--(gap2);
    \draw[arr] (gap2)--(exe); \draw[arr] (exe)--(rm); \draw[arr] (rm)--(log);
    \node[lbl, right=1mm of gap1, text=orange!60!black, align=left] {真机：相机直取帧};
    \node[lbl, right=1mm of gap2, text=orange!60!black, align=left] {真机：直调执行器};
  \end{tikzpicture}
  \caption{真机线主链与两个注入缝隙（橙色虚线）。缝两侧口径固定，故任务核心与执行侧可分别开发、
  失败可分别定位；两处缝隙正是真机相对仿真唯一改变的地方。}
  \label{fig:arch}
\end{figure}

\subsection{一次取放拆成六段：把动作序列固定，把“在哪儿”参数化}

“取一个物体并放好”是一个持续数秒的长任务，服务端把它切成固定的\textbf{六个阶段}，
每进入一段先发进度，再交给执行者 \texttt{PickExecutor}。
这个抽象是相对实验二最重要的一步泛化：实验二的直驱程序是“一个固定取物点 $\to$ 固定放置点”
的循环，而本实验有 6 格 $\times$ 2 料盒；若直接套用就会退化成“每格写死一个周期”。
真正会变的只是“\textbf{目标和料盒在哪儿}”，动作序列本身完全相同 ——
于是把序列固定成六个方法、把位置压进参数（表~\ref{tab:stage}）。

\begin{table}[htbp]
  \centering
  \small
  \begin{tabular}{@{}c l p{5.4cm} p{4.4cm}@{}}
    \toprule
    段 & 语义 & EP 原语 & 该段失败时的原因 \\
    \midrule
    1 & 下探 & \texttt{goto}(格位姿) $\to$ \texttt{arm\_moveto}(抓取低点) & \texttt{unreachable\_target} \\
    2 & 夹取 & \texttt{gripper\_close()} & \texttt{exec\_error} \\
    3 & 抬离 & \texttt{arm\_moveto}(抬升高点) $+$ 判 \texttt{held} & \texttt{grasp\_failed} \\
    4 & 搬运 & \texttt{goto}(料盒位姿) & \texttt{unreachable\_target} \\
    5 & 松放 & \texttt{arm\_moveto}(放置低点) $\to$ \texttt{gripper\_open()} $+$ 判 \texttt{placed} & \texttt{dropped\_or\_}\allowbreak\texttt{missed\_place} \\
    6 & 退回 & \texttt{arm\_moveto}(抬升高点) $\to$ \texttt{goto}(归位点) & \texttt{exec\_error} \\
    \bottomrule
  \end{tabular}
  \caption{六段动作与 EP 原语的对应，以及每段失败时映射出的动作级原因
  （\texttt{real/ep\_backend.py} 的 \texttt{STAGE\_REASON} 为唯一权威）。}
  \label{tab:stage}
\end{table}

这张表背后有一条我给自己定的规则，值得单独说：\textbf{只报这一段的代码真能做出来的判断，
不替谁宣布一个没人做过的判定。} 真机上六段\textbf{全都只能靠异常失败}（只有抬离段和松放段会
返回“否”，其余段没有可判的物理结果），所以这张表实际上是“\textbf{按段解释那次异常}”。
两个具体后果：

\begin{itemize}[leftmargin=1.6em]
  \item \textbf{夹取段报 \texttt{exec\_error} 而不是 \texttt{grasp\_failed}}：这一段的代码只是
        合一下夹爪，\textbf{从没判过“夹住没有”}（判定在下一段抬离）；而
        \texttt{grasp\_failed} 在原因枚举里的定义恰恰是“抬离后判定没夹起”。填它等于
        \textbf{凭空宣布一个没做过的判定}——现场查日志的人会以为系统真的判过。
  \item \textbf{退回段报 \texttt{exec\_error} 而不是 \texttt{unreachable\_target}}：仿真线这里
        用的是“不可达”，因为那个结论\textbf{来自规划器}；EP \textbf{没有规划层}，
        没有任何东西产出过“不可达”。报驱动硬错才是如实描述。
\end{itemize}

这也解释了为什么真机线与仿真线在\textbf{退回段}上的原因刻意保持不同：不是没对齐，
而是\textbf{两个平台上“发生了什么”本来就不是同一件事}。

\begin{figure}[htbp]
  \centering
  \begin{tikzpicture}[node distance=2mm]
    \node[proc, fill=blue!10] (s1) {\scriptsize 1\\下探\\\texttt{goto+moveto}};
    \node[proc, fill=blue!10, right=of s1] (s2) {\scriptsize 2\\夹取\\\texttt{close}};
    \node[proc, fill=blue!10, right=of s2] (s3) {\scriptsize 3\\抬离\\\texttt{moveto}$+$判};
    \node[proc, fill=blue!10, right=of s3] (s4) {\scriptsize 4\\搬运\\\texttt{goto}};
    \node[proc, fill=blue!10, right=of s4] (s5) {\scriptsize 5\\松放\\\texttt{open}$+$判};
    \node[proc, fill=blue!10, right=of s5] (s6) {\scriptsize 6\\退回\\\texttt{moveto+goto}};
    \node[hw, fill=red!10, below=7mm of s3] (f) {\scriptsize 任一段失败 $\Rightarrow$
      \textbf{立刻收尾}：开爪 $\to$ 抬到高 $\to$ 回归位点};
    \draw[arr] (s1)--(s2); \draw[arr] (s2)--(s3); \draw[arr] (s3)--(s4);
    \draw[arr] (s4)--(s5); \draw[arr] (s5)--(s6);
    \draw[arr, dashed, red!70] (f.north) -- (s3.south);
  \end{tikzpicture}
  \caption{六段动作链。与仿真线共用同一套段号；差别仅在原语是 EP SDK 而非 ROS 规划器。
  失败段号 $\to$ 原因 $\to$ 收尾这条路径见 §3.5。}
  \label{fig:stages}
\end{figure}

\subsection{EP 的位姿按编号查，不与桌面坐标换算（真机侧最关键的一处设计）}

仿真线里，“格号 $\to$ 坐标”用的是\textbf{桌面坐标系}——相机标定出来的那一套。真机上我
\textbf{没有}沿用这套，而是让每个格号/料盒号去位姿表里查一个\textbf{底盘里程计位姿}：

\begin{itemize}[leftmargin=1.6em]
  \item \textbf{为什么不换算：}EP 底盘的位姿是里程计读数，以\textbf{上一次上电的位置}为原点；
        相机标定出的桌面系以\textbf{桌面某一点}为原点。两者之间没有已知的对应关系（相机装在哪、
        盘停在哪都不固定），而\textbf{建立这个对应关系本身没有必要}——接口只给格号与类别。
  \item \textbf{于是：}两套坐标各管各的，永远不用互相换算；换格子、换料盒、挪相机都只改各自那份
        配置，\textbf{不动代码、不动接口}。“换格不改接口”这句话在真机侧的收益就落在这里。
  \item \textbf{代价（如实记录）：}底盘是\textbf{开环定位}，里程计会漂移，且\textbf{每次上电原点重置}。
        所以每次换场地或重新上电，位姿表必须重标——这一点被写进了启动前检查（§3.4）。
\end{itemize}

位姿表 \texttt{config\_real/ep\_waypoints.json} 的内容即为此：6 个网格位姿 + 2 个料盒位姿 +
1 个归位点（\texttt{goto} 的目标），臂的两档高度（\texttt{grab\_low\_mm} 抓取低点 /
\texttt{lift\_high\_mm} 抬升高点），以及夹爪的开合档位与动作等待时间。

\subsection{未标定拒跑闸门：把“标没标过”变成可执行的检查（个人认为本线最有价值的一步）}

\subsubsection*{问题的由来（实验二的教训）}

实验二栽过一次：配置文件里几何如果填的是\textbf{编出来的占位值}，照样能通过项目里全部的结构
自检——因为那些检查校验的是“三份配置之间\textbf{一致不一致}”，系统里\textbf{根本没有
“是否标定过”这个概念}。在仿真里这无所谓（坐标错了，屏幕上看得见）；在真机上，
那是\textbf{一条真机械臂开到一个没人量过的位置}。

\subsubsection*{做法}

\begin{enumerate}[leftmargin=1.8em]
  \item \textbf{未标定项一律写 \texttt{null}}——让“没标”在数据上就是一个\textbf{编不像}的值，
        而不是一个看着合理的数。占位项另外留 \texttt{\_note} 标记，闸门会扫描这些标记做兜底。
  \item \textbf{四份标定记录分开署名}（表~\ref{tab:calib}）——这才知道\textbf{卡在谁那里}，
        而不是笼统地报“配置没填完”。
  \item \textbf{闸门前置到“连机器人之前”}：真机入口 \texttt{run\_real.py} 启动的第一件事就是
        过闸门，未通过则直接退出，\textbf{相机不建、机器人不连}。
  \item \textbf{不给“跳过闸门”的开关}：想跳过只能去改配置里的标定记录，而那是\textbf{留痕}的。
\end{enumerate}

\begin{table}[htbp]
  \centering
  \small
  \begin{tabular}{@{}l l p{6.4cm}@{}}
    \toprule
    记录位置 & 负责人 & 内容 \\
    \midrule
    \texttt{task.json} $\to$ \texttt{camera.calibration} & 组员乙 & 相机标定（标定模型 + 数值 + 反投影误差） \\
    \texttt{grid\_cells.json} $\to$ \texttt{calibration} & 组长（现场量） & 取物网格几何（卷尺量出每格两角坐标） \\
    \texttt{bins.json} $\to$ \texttt{calibration} & 组长 & 分类料盒中心 \\
    \texttt{ep\_waypoints.json} $\to$ \texttt{calibration} & 组长（本人） & 8 个底盘位姿 + 臂两档 + 夹爪档位 \\
    \bottomrule
  \end{tabular}
  \caption{四份标定记录各自署名。任一份未填，真机入口拒绝启动。}
  \label{tab:calib}
\end{table}

\subsubsection*{为什么做成硬门而不是警告}

“未标定”这件事的风险不对称：误判“已标定”的代价是\textbf{真机撞桌或空抓}，而误判“未标定”的
代价只是\textbf{多跑一次命令}。风险不对称时，默认拒绝比默认放行合理。

\subsection{失败即收尾：段号 \texorpdfstring{$\to$}{->} 原因 \texorpdfstring{$\to$}{->} 立刻回到已知状态}

EP 是\textbf{开环}的：动作发出去之后，机械臂实际停在哪儿并没有可靠反馈。因此失败处置不能只
“记一笔往下走”——那样下一个物体开始时就等于\textbf{拖着一条位置未知的臂横扫桌面}。
本线的处理是三步（图~\ref{fig:fail}）：

\begin{enumerate}[leftmargin=1.8em]
  \item \textbf{段号定因：}失败发生在第几段，就决定了动作级原因是哪个（表~\ref{tab:stage}），
        日志里能直接看出“是没夹起”还是“够不着”；
  \item \textbf{立刻收尾 \texttt{safe\_reset}：}开爪 $\to$ 抬到高点 $\to$ 回归位点，
        \textbf{尽力而为、不再抛错}（反正已经在失败路径上）；
  \item \textbf{每次取放前先归位 \texttt{prepare()}：}上一次若是被人工中止或收尾也失败了，
        臂可能停在低处——所以每次动作前先抬到高位并开爪。这一步若失败\textbf{不吞}，
        直接停线。
\end{enumerate}

\begin{figure}[htbp]
  \centering
  \begin{tikzpicture}[node distance=4mm and 7mm]
    \node[proc] (base) {六段执行中};
    \node[dec, right=of base] (q) {某段失败？};
    \node[proc, right=of q, fill=orange!12] (map) {按\textbf{段号}定因\\（动作级原因）};
    \node[hw, right=of map, fill=red!8] (rst) {\texttt{safe\_reset}\\开爪 $\to$ 抬到高\\$\to$ 回归位点};
    \node[proc, below=6mm of base, fill=green!10] (nx) {下一物体：\\先 \texttt{prepare()} 归位};
    \draw[arr] (base)--(q);
    \draw[arr] (q)--node[above,lbl]{是}(map);
    \draw[arr] (map)--(rst);
    \draw[arr] (rst.south) |- (nx.east);
    \draw[arr] (q.south) |- (nx.east);
  \end{tikzpicture}
  \caption{失败路径：段号定因 $\to$ 立刻收尾 $\to$ 下一次动作前再归位。
  三级递进都是为了同一件事——\textbf{不让机械臂带着未知位姿继续工作}。}
  \label{fig:fail}
\end{figure}

\subsection{日志与验收口径：让“人工只做了该做的”也能被日志证明}

真机的日志落在 \texttt{exp3\_logs\_real/run\_<时间戳>/}，其中 \texttt{records.jsonl} 与
\texttt{result.json} 的 schema \textbf{与仿真线完全相同}（可直接比对两线结果），
另外多一份 \texttt{real\_run.json} 记录真机特有的上下文：是否演练、四份标定记录、
动作链步数、以及\textbf{人工确认的次数}。

最后一项是刻意留的凭据。EP 没有任何物体位姿感知，“夹住没有 / 放正没有”只能靠人确认；
而 docx §六 要求“\textbf{正常任务过程中不得进行人工类别判断和动作触发}”。两者的边界是：

\begin{itemize}[leftmargin=1.6em]
  \item \textbf{人不参与}：判类别（由识别模型给）、选目标与触发动作（由状态机给）；
  \item \textbf{人参与}：只回答“夹住了吗 / 放正了吗”这两个\textbf{关于物理姿态}的问题。
\end{itemize}

为了这条边界可核查，每条物体记录的 \texttt{detail} 字段会写明该次判断的\textbf{来源}——
是“操作员确认夹住”还是演练口径的“未判假定成功”，二者在日志里\textbf{用词不同、可区分}；
\texttt{real\_run.json} 另记人工确认总次数（正常一轮 = 2 $\times$ 物体数：
每个物体问一次 HELD、一次 PLACED）。验收时“人到底参与了什么”可以由文件读出，
而不是靠回忆。

\section{项目中遇到的问题和解决方法}

\textbf{问题 1：EP 的臂只有“一个水平自由度”，桌面横向覆盖从哪来。}

\textit{定位：}实验二的抓取动作可以直接照搬，但\textbf{可达范围}不能想当然。EP 工程形态的臂
\texttt{arm.moveto(x\_mm, y\_mm)} 中，$x$ 是\textbf{向前伸出的距离}、$y$ 是\textbf{高度}——
臂实际只在\textbf{一个竖直平面}内运动；官方行程约 0.22\,m（水平）$\times$ 0.15\,m（垂直），
需要现场实测。

\textit{解决：}把“横向覆盖”明确划给\textbf{底盘}：每个网格对应一个\textbf{底盘停车位姿}，
车先开到位、再由臂在竖直平面内完成下探/抬升。于是执行侧需要标定的不是“一格一个关节角”，
而是“\textbf{一格一个底盘位姿} + 一组臂档位”。

\textit{个人思考：}这条的教训是\textbf{先量工作空间再定场景}。如果按“桌面铺满网格、让臂去够”
的思路先把网格间距定死，真机上会直接撞上行程墙；而后台的对策（底盘逐格跑）代价是
\textbf{标定点数与里程计漂移都翻倍}。所以网格间距应当由“臂的实际可达域”倒推，
而不是先画好网格再让机器人去凑。

\medskip
\textbf{问题 2：EP 没有物块感知，“抓没抓住”无法自动判——但验收又禁止人工干预。}

\textit{定位：}夹爪没有力/位反馈，相机在动作过程中也会被臂遮挡，所以“是否夹住”“是否放正”
\textbf{在硬件上无从自动判断}。可 docx §六 明确要求正常过程中不得人工判类、不得人工触发动作。
两者看起来是冲突的。

\textit{解决：}把“人工干预”拆成两类区别对待——\textbf{决策性}的（判类别、选目标、发动作）
全部由系统做；“\textbf{物理状态确认}”性的（夹住了吗、放正了吗）由人回答，且\textbf{人的回答
不会触发任何动作}，只用于把这一物体记为成功或失败。并把这个边界\textbf{写进日志}（§3.6），
使“人只做了确认”这件事可核查，而不是一句自述。

\textit{个人思考：}“禁止人工干预”这类要求，真正要禁止的往往是\textbf{人在回路里做判断、当传感器}；
而“人作为物理状态的确认者”与它是两件事。把这条边界显式写出来并留下日志证据，
比在答辩时口头解释清楚得多。

\medskip
\textbf{问题 3：配置里“编的占位值”能骗过全部结构检查。}

\textit{定位：}项目原有的配置自检校验的是三份配置之间的一致性，\textbf{没有“是否标定过”这个概念}。
于是随手编一组“看着合理”的坐标，自检会\textbf{全部通过}——在仿真里无所谓，在真机上就是
让机械臂去一个没人量过的位置。

\textit{解决：}见 §3.4 的四条：未标定项写 \texttt{null}（编不出像样的默认值）、占位项留标记、
四份记录\textbf{分开署名}、闸门\textbf{前置到连机器人之前}且不给跳过开关。

\textit{个人思考：}检查项“通过”不等于“前提成立”。一个只检查内部一致性的校验器，对
“数据本身是否来自真实测量”是\textbf{盲}的；而这类盲区恰恰只在真机上才暴露。
把关键前提变成\textbf{可执行的门}，比在流程文档里写一句“记得先标定”可靠。

\medskip
\textbf{问题 4：开环平台上的失败处置——失败之后机械臂停在哪儿没人知道。}

\textit{定位：}EP 的动作是开环的，失败时没有可靠的位姿反馈。若只记一笔就继续下一个物体，
下一次动作就可能在“臂还停在桌面上方低处”的状态下开始，导致\textbf{拖着臂横扫}。

\textit{解决：}三级防线（§3.5）：失败即 \texttt{safe\_reset}（开爪 $\to$ 抬到高 $\to$ 回归位点，
尽力而为）；每次取放前 \texttt{prepare()} 先归位（失败则\textbf{拒绝动作}、宁可停线）；
退回段自带一次重试——物已经放下了，不值得为一次读位抖动把整轮判失败，但连试两次都回不去
就\textbf{真抛错}停线。

\textit{个人思考：}在开环系统里，“\textbf{回到一个已知状态}”比“尽量把当前动作做完”重要得多。
把收尾做成默认路径（而不是异常处理里的附加项），意味着每一段代码都必须回答
“我失败之后，机器人处在什么状态、谁负责收拾”。

\medskip
\textbf{问题 5：SDK 的“动作完成”回调早于实际到位。}

\textit{定位：}EP 的原语调用返回\textbf{动作成功}的时刻，比位置推送更新的时刻\textbf{早约 0.5\,s}
（官方把结果放在成功位里，而“等它返回”并不等于“动作已完成”）。若按“命令返回了”就往下走，
机械臂会在\textbf{移动途中}就开始抓取或搬运。

\textit{解决：}不自己另写一份驱动，而是\textbf{直接复用实验二已经调好的那份 EP 驱动}
（\texttt{ep/drive/rm.py}）。它内部的 \texttt{arm\_moveto()} 已经带了“等位置推送追上、
且连续两次读数一致才算稳定”的等待逻辑，本线的执行器调用它就自动继承了这套处理。

\textit{个人思考：}这次复用是为了\textbf{避免两份驱动各自漂移}——同一份 SDK 的坑，抄第二遍
一定会抄丢一部分。代价是公开快照里没有那份驱动，因此快照中只能做\textbf{不连硬件的演练}；
这个取舍我认为是划算的。

\medskip
\textbf{问题 6：越界命令会被“安静地钳制”。}

\textit{定位：}当目标高度超出行程时，SDK 会把命令\textbf{钳制}到边界，臂“稳定地停在”一个
并非目标的读数上。复用的驱动在检测到这种偏差时只打一条\textbf{警告日志、不抛错}——
于是执行侧会返回“已到抓取位”，是一个\textbf{假成功}。

\textit{解决：}运行时它仍会被兜住：钳制导致夹爪没对准物体 $\to$ 人工确认 HELD 为否 $\to$
判 \texttt{grasp\_failed} $\to$ 重抓一次 $\to$ 连续失败达上限 $\to$ 保护性停止。
\textbf{行为是安全的}，但现象容易被误读成“抓取不稳定”。因此把这条\textbf{明确写进
真机文档的已知限制}：现场若出现连续 \texttt{grasp\_failed}，第一件事是去日志里找那条
“未到位”警告，检查是不是标定值超了行程，而不是去怀疑夹爪。

\textit{个人思考：}“不抛错但结果不对”是最难查的一类失败。对这种情形，
即便代码选择继续运行，也必须在\textbf{文档里留下路标}，否则下一次联调还要重新推一遍。

\subsection*{附：本人这条线的复现方式}

\begin{lstlisting}[numbers=none]
# 1) 离线全套单测（无需 ROS、无需真机、无需相机）
cd exp3 && python3 -m unittest discover -s tests -t . -v     # 119 tests OK

# 2) 真机标定：读数直接写回位姿表，不用手抄（步骤见 real/标定说明.md）
python3 real/calibrate.py --odom                     # 推着车到位，实时看 odom
python3 real/calibrate.py --record c1                # 记该格底盘位姿
python3 real/calibrate.py --record-arm grab 74 120   # 试臂档（读回差 >3mm 拒绝入表）
python3 real/calibrate.py --seal --by <名字> --venue <场地>

# 3) 标定进度自查（未标定项逐个点名，含负责人）
python3 config_check.py --config-dir config_real
python3 config_check.py --config-dir config_real --require-calibrated   # 闸门口径，必须 0 项

# 4) 演练：不连 EP、不动真机，把整条动作链打出来（标定填完后先跑这个）
python3 real/run_real.py --config-dir config_real --dry-run \
        --camera none --detector mock --assume-judge --no-pause

# 5) 真机首跑：每个物体动手前等一次回车（人盯；本轮记为 rehearsal）
python3 real/run_real.py --config-dir config_real --confirm-each \
        --camera cv2:0 --detector file --detector-file <检测器文件>

# 6) 真机记分轮：三个演练开关全关，HELD / PLACED 向操作员确认
python3 real/run_real.py --config-dir config_real \
        --camera cv2:0 --detector file --detector-file <检测器文件>
#    日志：exp3_logs_real/run_<时间戳>/{run.log, records.jsonl, result.json, real_run.json}
#    real_run.json 记 evidence_grade 与 outcome，验收时据此判这份结果能不能用
\end{lstlisting}

\vspace{2mm}
{\small 说明：上列第 1--3 项（119 个离线用例，含逐段失败注入、闸门拦截、标定工具的
钳制读数拦截、四类异常的端到端复跑）在当前仓库中均可直接执行并已有结果。
\textbf{第 4--6 项中涉及 EP 硬件的部分属阶段二待完成}——本报告所述真机侧的设计与失败处置，
凡未标注“现场实测”者，均为\textbf{离线（假机械臂）验证}的结论，真机实测结果以现场记录与
验收日志为准。相关代码已整理在 \texttt{exp3/real/}（执行后端、标定工具、入口、相机与检测器）
与 \texttt{exp3/config\_real/}（真机配置与位姿表）中。}
\end{document}
